\chapter{Justificación}
\section{Justificación e importancia}
La implementación de estos sistemas es crucial para superar las limitaciones de los procesos manuales actuales en ambas áreas. Para la gestión de evidencias, se abordará la dificultad de acceso, organización y control, lo que se traducirá en:
\begin{itemize}
    \item Reducción significativa del tiempo de clasificación manual (hasta 70\% de ahorro).
    \item Eliminación de la duplicidad de archivos (abordando 35\% de discrepancias).
    \item Garantía del cumplimiento continuo de los estándares definidos por SINAES.
    \item Centralización efectiva de más de 500 documentos anuales con metadatos estructurados.
    \item Mejora de la seguridad y trazabilidad de la información sensible institucional.
\end{itemize}
Para el módulo de tiempos de jornada, el sistema representa una herramienta clave para la planificación eficiente de recursos humanos y académicos de la universidad, permitiendo tomar mejores decisiones sobre asignaciones docentes, ejecución de proyectos y administración del presupuesto anual de tiempos de jornada. La combinación de ambos módulos bajo el paraguas de "Gestión de Calidad" optimizará la toma de decisiones estratégicas y operativas en la UNA.
